\NSPage{151}%
uniformly in \NSi{NS-F-P151-001}{\mathsf C,t_1}; its real value is bounded above and below. Thus
\NSi{NS-F-P151-002}{\mathsf C E_{\mathrm r}=\sqrt{2X}\phi_{\mathrm r}} and its reciprocal have the asserted bounds away from
\NSi{NS-F-P151-003}{X=0}. Moreover
\NSFormula{NS-F-P151-004}{display}{B.24}
\begin{equation}
  \Pi_{\mathrm r}(X)-\Pi_0=\mathsf C^{-2}\int_0^X\phi_{\mathrm r}(X',\eta)^2\,dX',
  \qquad
  \sup_{X\leq X_i}\lVert\Pi_{\mathrm r}-\Pi_0\rVert_{C_\eta^k}\leq C_k\mathsf C^{-2}.
\tag{B.24}\label{eq:B.24}
\end{equation}
The regular axis factor makes this integral finite at zero.

For the source sign, its leading gradient contribution is
\NSi{NS-F-P151-005}{-H_*\widetilde\xi_0=L\Lambda\chi}. The bound
\NSi{NS-F-P151-006}{|\widetilde\xi_0|\leq L\Lambda\sigma_*^{-1}\sqrt\chi} shows that replacing
\NSi{NS-F-P151-007}{H_*} by \NSi{NS-F-P151-008}{H_{c,\mathrm r}} introduces an error bounded by
\NSi{NS-F-P151-009}{C_{\sigma_*}\sqrt\chi+o_{t_1}(1)}, after
\NSi{NS-F-P151-010}{\Lambda,\mathsf C} are fixed. The logarithmic slope of the axis profile and the remaining gradient give
\NSi{NS-F-P151-011}{C\chi+C/\Lambda}. Absorb these errors into a small fraction of
\NSi{NS-F-P151-012}{L\Lambda\chi}, using
\NSi{NS-F-P151-013}{C_{\sigma_*}\sqrt\chi\leq\varepsilon\Lambda L\chi+C_{\varepsilon,\sigma_*}/\Lambda}.
The constant contribution retains the slack from \NSi{NS-F-P151-014}{-W_*>2.8}, yielding the first inequality
in \eqref{eq:B.23}. Since the axis profile satisfies \NSi{NS-F-P151-015}{p_1=a>0}, its logarithmic
\NSi{NS-F-P151-016}{\phi} slope is nonpositive; cutting this slope to zero gives
\NSi{NS-F-P151-017}{l_{\mathrm r}\leq1}.

The axial source satisfies
\NSi{NS-F-P151-018}{S_{n,\mathrm r}=Z_*+O(\Lambda^{-1})+o_{t_1}(1)+O(\mathsf C^{-2})}. Here
\NSi{NS-F-P151-019}{\mathcal D_XU_{\mathrm{nat}}=O(\Lambda^{-1})} on the short cutoff, and \eqref{eq:B.24}
controls pressure and its needed parameter derivative on the entire finite interval. The equations for
\NSi{NS-F-P151-020}{p_{1,\mathrm r},n_{s,\mathrm r}} are
\NSFormula{NS-F-P151-021}{display}{B.25}
\begin{equation}
  \mathcal D_Xp_{1,\mathrm r}=\frac{XS_{q,\mathrm r}}{L}-l_{\mathrm r}p_{1,\mathrm r},
  \qquad
  \mathcal D_Xn_{s,\mathrm r}+n_{s,\mathrm r}=\frac{S_{n,\mathrm r}}{L}.
\tag{B.25}\label{eq:B.25}
\end{equation}
Their regular initial values and bounded coefficients give the claimed parameter bounds. On
\NSi{NS-F-P151-022}{\chi\leq.99}, the axis alternative
\NSi{NS-F-P151-023}{|Z_*|>\delta_*} therefore preserves the sign and a positive lower bound of
\NSi{NS-F-P151-024}{|n_{s,\mathrm r}|}.

The first equation in \eqref{eq:B.25} preserves positivity and gives both comparisons
\NSFormula{NS-F-P151-025}{display}{none}
\[
\begin{aligned}
  p_{1,\mathrm r}(y)&\geq p_{1,\mathrm r}(0)e^{-y}+1.88\chi(e^y-e^{-y}),\\
  p_{1,\mathrm r}(X)&\geq p_{1,\mathrm r}(X_0)\frac{X_0}{X}+1.2\left(X-\frac{X_0^2}{X}\right).
\end{aligned}
\]
On \NSi{NS-F-P151-026}{\chi>.99}, the first right-hand side remains above
\NSi{NS-F-P151-027}{2.3}, since its initial value can be replaced by
\NSi{NS-F-P151-028}{2.3} and \NSi{NS-F-P151-029}{1.88\chi>2.3/2}. On the complement,
\NSi{NS-F-P151-030}{p_{2,\mathrm r}=Xn_{s,\mathrm r}/E_{\mathrm r}} grows uniformly in absolute value with
\NSi{NS-F-P151-031}{\mathsf C}, while \NSi{NS-F-P151-032}{p_{1,\mathrm r}} remains positive and bounded. This proves
\NSi{NS-F-P151-033}{v_{\mathrm r}>2+c}. The second comparison gives
\NSi{NS-F-P151-034}{p_{1,\mathrm r}>3} from \NSi{NS-F-P151-035}{X=100} onward.
\end{proof}

\subsection{Continuation to nonzero leading residual stress.}\label{sec:B.6}
In this subsection, we use the reference bounds of Lemma~\ref{lem:B.4} to reduce the shear
\NSi{NS-F-P151-036}{\mathfrak s} while keeping the integrated residual coefficients
\NSi{NS-F-P151-037}{p_s} close to their reference values as in \eqref{eq:B.28}. The resulting stress
\NSi{NS-F-P151-038}{\mathcal T_0=F(p_s-\mathfrak s)} of \eqref{eq:4.11} must become nonzero smoothly at the end of the axis region,
satisfy the admissible stress cone on a first collar, and retain the relaxed cone condition along the rest of the
continuation (Proposition~\ref{prop:B.5}).

\begin{proposition}\label{prop:B.5}
The analytic axis profile can be continued to \NSi{NS-F-P151-039}{X_i} with positive
\NSi{NS-F-P151-040}{E}, remaining unchanged for \NSi{NS-F-P151-041}{X\leq X_0}. Its stress has the factorization
\NSi{NS-F-P151-042}{\mathcal T_0=e_aB_0}, where the scalar factor \NSi{NS-F-P151-043}{e_a} vanishes to infinite order at
\NSi{NS-F-P151-044}{X_0} and the smooth vector coefficient \NSi{NS-F-P151-045}{B_0} is nonzero there. The admissible stress
cone holds on a first collar, and the strict relaxed cone condition holds on the rest of
\NSi{NS-F-P151-046}{(X_0,X_i]}. At \NSi{NS-F-P151-047}{X_i},
\NSi{NS-F-P151-048}{a=.8}, \NSi{NS-F-P151-049}{\mathcal D_XU=0}, and \NSi{NS-F-P151-050}{p_1>2}.
\end{proposition}
\begin{proof}
We reduce the reference shear on a short interval while \NSi{NS-F-P151-051}{p_s} changes by a controlled
small error. The resulting residual stress is nonzero and satisfies the cone conditions stated
above. We then verify its factorization by a function vanishing to infinite order and complete
the continuation with small shear.

\NSPage{152}%
Fix a sufficiently large finite \NSi{NS-F-P152-001}{\mathsf C}. Choose
\NSi{NS-F-P152-002}{0<t_*\leq\bar t} sufficiently small that the bounds of Lemma~\ref{lem:B.4}
hold uniformly for every reference width \NSi{NS-F-P152-003}{0<t_1\leq t_*}. The parameter bounds of
Lemma~\ref{lem:B.4} and its positive lower comparison for \NSi{NS-F-P152-004}{p_{1,\mathrm r}} give one finite upper bound
\NSi{NS-F-P152-005}{V_{\max}} for \NSi{NS-F-P152-006}{v_{\mathrm r}} over this entire reference family. It may depend on
\NSi{NS-F-P152-007}{\mathsf C}, but not on \NSi{NS-F-P152-008}{t_1}. Let
\NSi{NS-F-P152-009}{0<\kappa_0<1/2}, to be chosen using that bound, and on
\NSi{NS-F-P152-010}{0<y<t_1} set
\NSFormula{NS-F-P152-011}{display}{B.26}
\begin{equation}
  e_a=(1-\kappa_0)\sigma(y/t_1),\quad
  \kappa=1-e_a,\quad
  a=\kappa p_{1,\mathrm r},\quad
  \mathcal D_XU=-\frac{\kappa Xn_{s,\mathrm r}}{2},\quad
  \partial_y\log\phi=-\frac{\kappa p_{1,\mathrm r}}{2}.
\tag{B.26}\label{eq:B.26}
\end{equation}
The reference agrees with the analytic axis profile on this interval. Subtracting its equations
gives exactly
\NSFormula{NS-F-P152-012}{display}{B.27}
\begin{equation}
  \partial_y(\log\phi-\log\phi_{\mathrm r})=\frac{e_ap_{1,\mathrm r}}{2},
  \qquad
  \partial_y(U-U_{\mathrm r})=\frac{e_aXn_{s,\mathrm r}}{2}.
\tag{B.27}\label{eq:B.27}
\end{equation}
Since \NSi{NS-F-P152-013}{e_a} increases, each fixed parameter derivative of these field differences is
\NSi{NS-F-P152-014}{O(ye_a)}. Constants in the activation comparisons may depend on the already fixed
\NSi{NS-F-P152-015}{\mathsf C,\Lambda}, but are uniform as \NSi{NS-F-P152-016}{t_1} and
\NSi{NS-F-P152-017}{\kappa_0} decrease within the admissible ranges. The reference family bounds
control the integrated coefficients independently of both widths. The radial moment identities
\eqref{eq:4.16}, using one additional parameter derivative, imply
\NSFormula{NS-F-P152-018}{display}{B.28}
\begin{equation}
  p_s-p_{s,\mathrm r}=O(ye_a),
  \qquad
  E_{\mathrm r}/E=1+O(ye_a).
\tag{B.28}\label{eq:B.28}
\end{equation}
The comparison holds for field values and parameter derivatives; radial derivatives of the
narrow cutoffs may be large.

The actual shear is \NSi{NS-F-P152-019}{\kappa(p_{1,\mathrm r},p_{2,\mathrm r}E_{\mathrm r}/E)}. In particular the common
\NSi{NS-F-P152-020}{\kappa} cancels from the ratio \NSi{NS-F-P152-021}{t_s=-b_s/a} before it is estimated.
The definitions of \NSi{NS-F-P152-022}{P_{\mathrm c},J_{\mathrm c},v_s} yield
\NSFormula{NS-F-P152-023}{display}{B.29}
\begin{equation}
  t_s=\frac{p_{2,\mathrm r}E_{\mathrm r}}{p_{1,\mathrm r}E},\qquad
  v_s=\kappa v_{\mathrm r}+O(ye_a),\qquad
  P_{\mathrm c}=v_{\mathrm r}+O(ye_a),\qquad
  J_{\mathrm c}=O(ye_a).
\tag{B.29}\label{eq:B.29}
\end{equation}
No inverse power of \NSi{NS-F-P152-024}{\kappa_0} enters these error constants. Taking
\NSi{NS-F-P152-025}{t_1} sufficiently small gives
\NSi{NS-F-P152-026}{P_{\mathrm c}-v_s\geq ce_a}, \NSi{NS-F-P152-027}{P_{\mathrm c}>2+c}, and
\NSFormula{NS-F-P152-028}{display}{none}
\[
  (v_s-2)_+J_{\mathrm c}^2<2(P_{\mathrm c}-v_s)^2
  \qquad(0<y\leq t_1).
\]
Indeed the ratio of its left side to \NSi{NS-F-P152-029}{(P_{\mathrm c}-v_s)^2} is at most
\NSi{NS-F-P152-030}{Cy^2}. The quadratic test of Lemma~\ref{lem:4.5} proves the admissible stress cone condition
when \NSi{NS-F-P152-031}{v_s>2}; when \NSi{NS-F-P152-032}{v_s\leq2},
\NSi{NS-F-P152-033}{P_{\mathrm c}>2} gives the strict relaxed condition directly. Continuity from
\NSi{NS-F-P152-034}{v_{\mathrm r}(0)>2+c} gives one positive-width first collar on which
\NSi{NS-F-P152-035}{v_s>2+c/2}, uniformly in \NSi{NS-F-P152-036}{\eta} for the final fixed parameters.

To establish smoothness after division by the flat factor, write
\NSi{NS-F-P152-037}{e_a=e^{-t_1^2/y^2}g_a(y)} near zero, where
\NSi{NS-F-P152-038}{g_a} is smooth and positive. Lemma~\ref{lem:A.9}, applied with
\NSi{NS-F-P152-039}{c=t_1^2}, \NSi{NS-F-P152-040}{j=0}, and coefficient
\NSi{NS-F-P152-041}{g_a b}, gives for every smooth coefficient \NSi{NS-F-P152-042}{b}
\NSFormula{NS-F-P152-043}{display}{none}
\[
  \frac{1}{e_a(y)}\int_0^y e_a(u)b(u,\eta)\,du
  =y^3B(y,\eta),
  \qquad B\in C^\infty.
\]
Division by the positive smooth \NSi{NS-F-P152-044}{g_a(y)} preserves smoothness and bounds on mixed derivatives
of each fixed order. The differences in \eqref{eq:B.27} belong to
\NSi{NS-F-P152-045}{e_ay^3C^\infty}. Products, smooth compositions, and reciprocals of nonvanishing fields preserve this class. The moment and
pressure differences vanish before \NSi{NS-F-P152-046}{X_0}; their additional integration puts them in
\NSi{NS-F-P152-047}{e_ay^6C^\infty}. Substitution in \eqref{eq:4.16} consequently justifies \eqref{eq:B.28}
after division by \NSi{NS-F-P152-048}{e_a} in every fixed mixed derivative, and gives the exact factorization
\NSFormula{NS-F-P152-049}{display}{B.30}
\begin{equation}
  \mathcal T_0=e_aB_0(y,\eta),
  \qquad
  B_0(0,\eta)=F(X_0,\eta)p_{s,\mathrm r}(X_0,\eta)\neq0,
  \qquad
  B_0\in C^\infty.
\tag{B.30}\label{eq:B.30}
\end{equation}
\NSPage{153}%
At the edge the limiting shear is \NSi{NS-F-P153-001}{p_{s,\mathrm r}}. Therefore
\NSi{NS-F-P153-002}{B_0\cdot(-t_s,1)=0} and \NSi{NS-F-P153-003}{B_0\cdot(1,t_s)>0} there. At
the endpoint, \NSi{NS-F-P153-004}{B_0\cdot(1,t_s)>0} and
\NSi{NS-F-P153-005}{2(B_0\cdot(1,t_s))^2-(v_s-2)(B_0\cdot(-t_s,1))^2>0}. Both quantities
have positive lower bounds on a smaller collar. In particular, for every fixed radial/parameter
multi-index \NSi{NS-F-P153-006}{I},
\NSFormula{NS-F-P153-007}{display}{B.31}
\begin{equation}
  |\mathcal T_0|\geq ce^{-t_1^2/y^2},
  \qquad
  |\partial^I\mathcal T_0|\leq C_Ie^{-t_1^2/y^2}y^{-N_I}.
\tag{B.31}\label{eq:B.31}
\end{equation}
The derivatives in \eqref{eq:B.31} act in \NSi{NS-F-P153-008}{(y,\eta)}; fixed derivatives in
\NSi{NS-F-P153-009}{X} obey the same form because \NSi{NS-F-P153-010}{X_0>0}. The nonzero vector
\NSi{NS-F-P153-011}{B_0(0,\eta)} determines the limiting stress direction.

We now continue to \NSi{NS-F-P153-012}{X_i} while preserving the relaxed cone condition. After the first
transition keep \eqref{eq:B.26} with \NSi{NS-F-P153-013}{\kappa=\kappa_0} through the reference cutoff and constant-slope interval.
For fixed \NSi{NS-F-P153-014}{\mathsf C,\Lambda}, integration over \NSi{NS-F-P153-015}{[X_0,X_{\mathrm b}]} and \eqref{eq:4.16} give
\NSi{NS-F-P153-016}{O(t_1+\kappa_0)} differences of field values, logarithms, and the required derivatives with respect to
\NSi{NS-F-P153-017}{\eta}. Choose \NSi{NS-F-P153-018}{\kappa_0} small compared to
\NSi{NS-F-P153-019}{V_{\max}} and the uniform family comparison constants, then choose
\NSi{NS-F-P153-020}{t_1} small. This order does not require a reference width to be fixed before
\NSi{NS-F-P153-021}{\kappa_0}. Formula \eqref{eq:B.29} now gives
\NSi{NS-F-P153-022}{P_{\mathrm c}>2+c/2} and \NSi{NS-F-P153-023}{v_s<1}. The strict relaxed cone condition follows,
regardless of the size of \NSi{NS-F-P153-024}{J_{\mathrm c}}.

At \NSi{NS-F-P153-025}{X_{\mathrm b}}, multiply the axial prescription for \NSi{NS-F-P153-026}{\mathcal D_XU} by a smooth factor
\NSi{NS-F-P153-027}{\beta} decreasing from one to zero in a short log interval. During this transition,
\NSi{NS-F-P153-028}{P_{\mathrm c}=p_{1,\mathrm r}+\beta p_{2,\mathrm r}^2/p_{1,\mathrm r}+o(1)>2}, while
\NSi{NS-F-P153-029}{v_s<1}. Next interpolate \NSi{NS-F-P153-030}{a} convexly from
\NSi{NS-F-P153-031}{\kappa_0p_{1,\mathrm r}} to \NSi{NS-F-P153-032}{.8}, keeping
\NSi{NS-F-P153-033}{\mathcal D_XU=0}. Arrange first that
\NSi{NS-F-P153-034}{\kappa_0\sup p_{1,\mathrm r}<.8}. The second transition has
\NSi{NS-F-P153-035}{v_s=a\leq.8}, and \NSi{NS-F-P153-036}{P_{\mathrm c}=p_1>2} by continuity from
the positive margin established at \NSi{NS-F-P153-037}{X_{\mathrm b}}. Both transitions fit strictly before
\NSi{NS-F-P153-038}{X_i}. On the final constant-slope interval
\NSi{NS-F-P153-039}{a=.8}, \NSi{NS-F-P153-040}{l=.6}, and the same gradient absorption as in Lemma~\ref{lem:B.4}
gives \NSi{NS-F-P153-041}{S_q>1}: its constant contribution is
\NSi{NS-F-P153-042}{.6(-W_*)>1.68}, while its leading gradient is
\NSi{NS-F-P153-043}{L\Lambda\chi}. At a hypothetical downward crossing
\NSi{NS-F-P153-044}{p_1=2},
\NSi{NS-F-P153-045}{\mathcal D_Xp_1=XS_q/L-.6p_1>X-1.2>0}, a contradiction. This completes the continuation.
\end{proof}

The continuation remains analytic in \NSi{NS-F-P153-046}{\eta} on the first collar. We reserve a smaller rectangle
there, so that later modifications preserve this regularity as well as the inner stress factorization.

\begin{corollary}\label{cor:B.6}
There is \NSi{NS-F-P153-047}{t_c>0}, with \NSi{NS-F-P153-048}{t_c<t_1} and
\NSi{NS-F-P153-049}{4e^{t_c}<4.1}, such that the functions
\NSi{NS-F-P153-050}{E/\sqrt{2X},U,V_0/X,\Pi} on
\NSi{NS-F-P153-051}{[0,X_0e^{t_c}]\times[-1,1]} are smooth in
\NSi{NS-F-P153-052}{(X,\eta)} through \NSi{NS-F-P153-053}{X=0} and analytic in
\NSi{NS-F-P153-054}{\eta} on one complex neighborhood. Every fixed radial derivative of these four functions has that same
parameter neighborhood and finite bounds on a smaller neighborhood. On
\NSi{NS-F-P153-055}{X_0<X\leq X_0e^{t_c}}, the admissible stress cone condition and the factorization
\eqref{eq:B.30} hold. All subsequent profile modifications are supported strictly to the right of
\NSi{NS-F-P153-056}{X_0e^{t_c}}.
\end{corollary}
\begin{proof}
Choose a compact complex neighborhood inside that of Proposition~\ref{prop:B.2}, small enough that the positive function
\NSi{NS-F-P153-057}{\Phi} of the analytic axis profile has no zeros. The formulas for
\NSi{NS-F-P153-058}{Q_s,N_s} computed from the reference profile, and \eqref{eq:B.26}, involve only analytic parameter coefficients,
their parameter derivatives, and forward integrals, multiplied by cutoffs in \NSi{NS-F-P153-059}{y} alone. Their
denominators are nonzero on a smaller neighborhood. Exponentiation gives nonvanishing
\NSi{NS-F-P153-060}{\phi}, so one neighborhood works throughout the fixed first collar. Radial differentiation changes
the cutoff bounds but not the parameter domain. Compactness and Proposition~\ref{prop:B.5} give an inner collar on which the admissible stress cone condition holds; choose
\NSi{NS-F-P153-061}{t_c} strictly inside it and reserve the resulting rectangle from all later modifications. Neither
\NSi{NS-F-P153-062}{t_c} nor radial derivative bounds are asserted uniform as \NSi{NS-F-P153-063}{\mathsf C} grows and the transition widths shrink. The
radial moment and pressure integrals are forward, so changes supported later leave this
rectangle unchanged.
\end{proof}

\NSPage{154}%
\subsection{A radial transition length chosen before the amplitude.}\label{sec:B.7}
In this subsection, we construct the remaining transition \eqref{eq:B.34} to
\NSi{NS-F-P154-001}{E=P_*f(X/X_R)^{1/10}}, with
\NSi{NS-F-P154-002}{f=(1+\eta^2)^{-1}}, before correcting its radial moment functions in the next subsection
(Proposition~\ref{prop:B.8}). We first bound its logarithmic radial length
\NSi{NS-F-P154-003}{T_{\mathrm{sh}}} independently of the amplitude
\NSi{NS-F-P154-004}{\mathsf C} chosen later, using Lemma~\ref{lem:B.7} and \eqref{eq:B.33}. Increasing that amplitude can then move the matching radius
\NSi{NS-F-P154-005}{X_R} outward without lengthening this transition, making the normalized inner moment discrepancies
small (Equations~\eqref{eq:B.38} and \eqref{eq:B.39}). The estimates below concern parameter derivatives;
bounds for derivatives with respect to \NSi{NS-F-P154-006}{X} need not be uniform in the radial transition widths.
In the bounds below, \NSi{NS-F-P154-007}{\kappa_0\in(0,1/2)} is the parameter controlling the shear reduction in \eqref{eq:B.26}.

\begin{lemma}\label{lem:B.7}
Fix the data for the outer profile and the axis, \NSi{NS-F-P154-008}{\Lambda}, and a finite parameter order
\NSi{NS-F-P154-009}{k}. Put \NSi{NS-F-P154-010}{\ell_i=\log(\mathsf C E(X_i,\eta))} and
\NSi{NS-F-P154-011}{G_i=U(X_i,\eta)}. Let \NSi{NS-F-P154-012}{\omega_{\mathrm{fin}}} be the sum of the logarithmic radial lengths of the
final axial cutoff and the interpolation of \NSi{NS-F-P154-013}{a} to
\NSi{NS-F-P154-014}{.8} in the proof of Proposition~\ref{prop:B.5}. The continuations
constructed there satisfy
\NSFormula{NS-F-P154-015}{display}{B.32}
\begin{equation}
\begin{aligned}
  \lVert\ell_i\rVert_{C_\eta^k}&\leq\mathcal B_k,\\
  \lVert G_i-U_*\rVert_{C_\eta^k}&\leq\frac{C_k^{\mathrm{nat}}}{\Lambda}
    +C_k^{\mathrm{join}}(\Lambda)(t_1+\kappa_0+\omega_{\mathrm{fin}}).
\end{aligned}
\tag{B.32}\label{eq:B.32}
\end{equation}
The constant \NSi{NS-F-P154-016}{C_k^{\mathrm{nat}}} is independent of both
\NSi{NS-F-P154-017}{\Lambda} and sufficiently large \NSi{NS-F-P154-018}{\mathsf C}, with the preceding axis data fixed.
The bounds \NSi{NS-F-P154-019}{\mathcal B_k} and \NSi{NS-F-P154-020}{C_k^{\mathrm{join}}(\Lambda)} may depend on
\NSi{NS-F-P154-021}{\Lambda}, but are independent of sufficiently large \NSi{NS-F-P154-022}{\mathsf C} and of
smaller radial transition widths.
\end{lemma}
\begin{proof}
Write \NSi{NS-F-P154-023}{L_0=\log(X_i/X_0)}. Let \NSi{NS-F-P154-024}{M_k,N_k} be the amplitude-independent bounds for
\NSi{NS-F-P154-025}{p_{1,\mathrm r},n_{s,\mathrm r}} from Lemma~\ref{lem:B.4}. Before the final interpolation of
\NSi{NS-F-P154-026}{a} to \NSi{NS-F-P154-027}{.8}, the actual logarithmic slope is
\NSi{NS-F-P154-028}{-\kappa p_{1,\mathrm r}/2}, with \NSi{NS-F-P154-029}{0<\kappa\leq1}; during that interpolation it varies convexly to
\NSi{NS-F-P154-030}{-.4}, after which it remains constant. Consequently
\NSFormula{NS-F-P154-031}{display}{none}
\[
  \lVert\log\phi(X_i)\rVert_{C_\eta^k}
  \leq\lVert\log\phi_*+\log\Phi(4,\mathord\cdot)\rVert_{C_\eta^k}
  +\tfrac12(M_k+.8)L_0.
\]
Adding \NSi{NS-F-P154-032}{\tfrac12\log(2X_i)} proves the first assertion, using \eqref{eq:B.13}. The first axial transition changes
\NSi{NS-F-P154-033}{U} by at most \NSi{NS-F-P154-034}{X_0e^{\bar t}N_kt_1/2}; the remainder contributes at most
\NSi{NS-F-P154-035}{X_iN_k\kappa_0/2}, with the same bound for its transition to zero. The approximation error for the analytic axis profile is
\NSi{NS-F-P154-036}{C_k^{\mathrm{nat}}/\Lambda}, with the constant independent of
\NSi{NS-F-P154-037}{\Lambda}, by \eqref{eq:B.13}. These estimates prove \eqref{eq:B.32}. Only integrals
of bounded cutoff values were used, so inverse powers of their widths do not occur.
\end{proof}

The bounds just proved let us choose the length of the transition independently of the
final amplitude. We choose
\NSFormula{NS-F-P154-038}{display}{B.33}
\begin{equation}
  T_{\mathrm{sh}}\geq20\lVert\sigma'\rVert_\infty
  \bigl(\mathcal B_0+\lVert\log f\rVert_\infty\bigr).
\tag{B.33}\label{eq:B.33}
\end{equation}
For \NSi{NS-F-P154-039}{y_i=\log(X/X_i)} we prescribe
\NSFormula{NS-F-P154-040}{display}{B.34}
\begin{equation}
  \log E=-\log\mathsf C+\frac{y_i}{10}
  +(1-\sigma(y_i/T_{\mathrm{sh}}))\ell_i
  +\sigma(y_i/T_{\mathrm{sh}})\log f,
  \qquad U=G_i.
\tag{B.34}\label{eq:B.34}
\end{equation}
Then \NSi{NS-F-P154-041}{l=.6+T_{\mathrm{sh}}^{-1}\sigma'(y_i/T_{\mathrm{sh}})(\log f-\ell_i)} lies in
\NSi{NS-F-P154-042}{[.55,.65]}, so \NSi{NS-F-P154-043}{a\in[.7,.9]} and
\NSi{NS-F-P154-044}{b_s=0}.

Along the transition \eqref{eq:B.34}, the source still satisfies \NSi{NS-F-P154-045}{S_q>1}. To verify its sign, make the
changes before the final \NSi{NS-F-P154-046}{.8} constant-slope interval small in the required
\NSi{NS-F-P154-047}{C_\eta^k} norms; that constant-slope interval has a parameter-independent slope. Thus
\NSi{NS-F-P154-048}{\ell_i'=\widetilde\xi_0+(\log\Phi(4,\mathord\cdot))_\eta+}%
\NSPage{155}%
\NSi{NS-F-P155-001}{o(1)} and
\NSi{NS-F-P155-002}{G_i-U_*=O(\Lambda^{-1})+o(1)}. The old gradient obeys
\NSFormula{NS-F-P155-003}{display}{none}
\[
\begin{aligned}
  -H_c\ell_i'&\geq L\Lambda\chi-C\chi-C_{\sigma_*}\sqrt\chi-C/\Lambda-o(1),\\
  -H_c(\log f)'&=\frac{2\eta H_c}{1+\eta^2}\geq-Cj_0^2-C/\Lambda-o(1).
\end{aligned}
\]
For the second line use \NSi{NS-F-P155-004}{\eta H_*=(D+4d)\eta^2+dj_0\eta} and complete the square. For the first line
use \eqref{eq:B.20}, which gives
\NSi{NS-F-P155-005}{H_*\partial_\eta\log\Phi=O(\chi)+O(\Lambda^{-1})}; multiplication of the
\NSi{NS-F-P155-006}{O(\Lambda^{-1})} field error by \NSi{NS-F-P155-007}{\widetilde\xi_0} introduces an error bounded by
\NSi{NS-F-P155-008}{C_{\sigma_*}\sqrt\chi}, which is absorbed as before. The two gradients are combined convexly in
\eqref{eq:B.34}. Since \NSi{NS-F-P155-009}{-W} remains close to
\NSi{NS-F-P155-010}{-W_*>2.8}, \NSi{NS-F-P155-011}{l\geq.55}, and
\NSi{NS-F-P155-012}{h,j_0} are small, \NSi{NS-F-P155-013}{S_q>1} follows. The barrier
\NSi{NS-F-P155-014}{\mathcal D_Xp_1=XS_q/L-lp_1>0} at \NSi{NS-F-P155-015}{p_1=2} preserves
\NSi{NS-F-P155-016}{p_1>2}. Thus this entire transition, and the subsequent constant-slope interval before
restoration, continue to satisfy the strict inequalities defining the relaxed cone condition.

\subsection{Matching the five radial moment functions exactly.}\label{sec:B.8}
The azimuthal profile \NSi{NS-F-P155-017}{E} now has the required reference form \eqref{eq:A.7}, and the axial profile
\NSi{NS-F-P155-018}{U=G_i} is close to its reference value \NSi{NS-F-P155-019}{4\eta} by \eqref{eq:B.32} and \eqref{eq:B.1}.
In this subsection, we match the five cumulative radial integrals
\NSi{NS-F-P155-020}{M,I,J,S,C_p} of \eqref{eq:4.15} to attach the outer profile. We first choose the radial scale
\NSi{NS-F-P155-021}{X_R} so that their normalized discrepancies are small, then correct them with five fixed bumps
while preserving the relaxed cone inequalities (Proposition~\ref{prop:B.8}). The common axis pressure datum
\NSi{NS-F-P155-022}{\Pi_0} and exact moment matching will then preserve every subsequent outer field by Lemma~\ref{lem:4.4}.

We set
\NSFormula{NS-F-P155-023}{display}{none}
\[
  X_R=X_i(\mathsf C P_*)^{10},
  \qquad
  x=X/X_R.
\]
We also set \NSi{NS-F-P155-024}{X_{\mathrm{sep}}=X_ie^{T_{\mathrm{sh}}}}, the endpoint of the transition
\eqref{eq:B.34}; it is fixed before the amplitude.

\begin{proposition}\label{prop:B.8}
For a sufficiently large amplitude \NSi{NS-F-P155-025}{\mathsf C}, followed by sufficiently small activation transitions,
the profile constructed in Proposition~\ref{prop:B.5} and continued by \eqref{eq:B.34} can be continued to
\NSi{NS-F-P155-026}{\log x=-5}, preserving the strict relaxed cone condition, so that its fields, pressure
\NSi{NS-F-P155-027}{\Pi}, and all five radial moment functions
\NSi{NS-F-P155-028}{M,I,J,S,C_p} agree exactly with those of the reference inner profile \eqref{eq:A.7}. It
can then follow the reference outer radial profile without changing its prescribed pressure datum or
the subsequent values of \NSi{NS-F-P155-029}{Q_s,N_s}.
\end{proposition}
\begin{proof}
We first normalize the radial moment map so that its inverse bounds do not depend
on the later large radius. We then show that the discrepancy accumulated at the inner
endpoint is small in those units and cancel it with five fixed bumps.

We introduce \NSi{NS-F-P155-030}{\widehat H=\sqrt{2x}E} and normalized radial moment functions by
\NSFormula{NS-F-P155-031}{display}{none}
\[
\begin{gathered}
  M=X_R\widehat M,\qquad
  I=X_R^{3/2}\widehat I,\qquad
  J=X_R^{3/2}\widehat J,\qquad
  S=X_R\widehat S,\\
  C_p=\displaystyle\int_0^x\frac{E(\xi,\eta)^2}{2\xi}\,d\xi,
  \qquad
  \Pi=\Pi_0+C_p.
\end{gathered}
\]
Here \NSi{NS-F-P155-032}{\widehat M,\widehat I,\widehat J,\widehat S} are the integrals of
\NSi{NS-F-P155-033}{U,\widehat H,U\widehat H,U^2-E^2/2}, respectively, with respect to
\NSi{NS-F-P155-034}{x}. Substitution in \eqref{eq:4.16} gives
\NSFormula{NS-F-P155-035}{display}{B.35}
\begin{equation}
\begin{aligned}
  W&=1-2D\eta\widehat M/x-d\widehat M_\eta/x,\\
  Q_s&=-W+\frac{(1-h)\widehat I-D\eta\widehat I_\eta-d\widehat J_\eta
                   +2(h-D)\eta\widehat J}{x\widehat H},\\
  N_s&=-WU+\frac{D(\widehat M-\eta\widehat M_\eta)+4h\eta\widehat S-d\widehat S_\eta}{x}
                   +4A\eta\Pi-d\Pi_\eta.
\end{aligned}
\tag{B.35}\label{eq:B.35}
\end{equation}
\NSPage{156}%
In particular \NSi{NS-F-P156-001}{X_R} has disappeared from the formulas for
\NSi{NS-F-P156-002}{Q_s,N_s}. On the fixed interval
\NSi{NS-F-P156-003}{\mathcal R=[e^{-8},e^{-5}]}, the ideal
\NSi{NS-F-P156-004}{E_0=P_*fx^{1/10}} and \NSi{NS-F-P156-005}{Q_{s,0}} have positive minima, and the ideal
\NSi{NS-F-P156-006}{N_{s,0}} is bounded; these are properties of the explicit formula \eqref{eq:A.25}. The map \eqref{eq:B.35} is Lipschitz
there in each prescribed finite parameter norm, with one further parameter derivative on its
inputs, and with constants depending only on data of the outer profile.

The same is true of the fixed restoration and moment correction operations. To make
the cone tolerance explicit, set
\NSi{NS-F-P156-007}{Q_{\min}=\min_{\mathcal R\times[-1,1]}Q_{s,0}>0},
\NSi{NS-F-P156-008}{e_*=\min E_0>0}, and
\NSi{NS-F-P156-009}{w_*=1+\max|N_{s,0}/(E_0Q_{s,0})|}. A tolerance determined by the exterior-profile data can ensure
throughout both operations that
\NSFormula{NS-F-P156-010}{display}{B.36}
\begin{equation}
  Q_s\geq Q_{\min}/2,\qquad
  E\geq e_*/2,\qquad
  |a-.8|\leq.1,\qquad
  \left|\frac{N_s}{EQ_s}\right|\leq w_*,\qquad
  |b_s|\leq\frac{.1}{1+w_*}.
\tag{B.36}\label{eq:B.36}
\end{equation}
It follows that \NSi{NS-F-P156-011}{v_s\leq.9+.01/.7<1}. Setting
\NSi{NS-F-P156-012}{G=Q_s-b_sN_s/(aE)}, we obtain
\NSFormula{NS-F-P156-013}{display}{B.37}
\begin{equation}
  G\geq\frac67Q_s\geq\frac{3Q_{\min}}7,
  \qquad
  P_{\mathrm c}=\frac{X_Rx}{L}G.
\tag{B.37}\label{eq:B.37}
\end{equation}
Thus \NSi{NS-F-P156-014}{P_{\mathrm c}>2} follows by increasing \NSi{NS-F-P156-015}{X_R} after fixing the tolerance. No moment correction
tolerance has to shrink with \NSi{NS-F-P156-016}{X_R}.

It remains to show that this fixed tolerance is achievable. In normalized coordinates, the
endpoint \NSi{NS-F-P156-017}{X_{\mathrm{sep}}} has coordinate
\NSFormula{NS-F-P156-018}{display}{B.38}
\begin{equation}
  x_{\mathrm{sep}}=X_{\mathrm{sep}}/X_R=e^{T_{\mathrm{sh}}}/(\mathsf C P_*)^{10}.
\tag{B.38}\label{eq:B.38}
\end{equation}
On \NSi{NS-F-P156-019}{[0,X_{\mathrm{sep}}]}, the previously established bounds for
\NSi{NS-F-P156-020}{\log(\mathsf C E(X_i,\mathord\cdot))} and
\NSi{NS-F-P156-021}{U(X_i,\mathord\cdot)}, together with \eqref{eq:B.34} bound
\NSi{NS-F-P156-022}{U} and \NSi{NS-F-P156-023}{\mathsf C E} in each required parameter norm; near zero
\NSi{NS-F-P156-024}{\mathsf C E=O(\sqrt X)}. Hence the contributions to the cumulative integrals from
\NSi{NS-F-P156-025}{0\leq x\leq x_{\mathrm{sep}}} satisfy
\NSFormula{NS-F-P156-026}{display}{B.39}
\begin{equation}
\begin{aligned}
  \lVert\widehat M\rVert_{C_\eta^k}+\lVert\widehat S\rVert_{C_\eta^k}
    &\leq C_kx_{\mathrm{sep}},\\
  \lVert\widehat I\rVert_{C_\eta^k}+\lVert\widehat J\rVert_{C_\eta^k}
    &\leq C_k\mathsf C^{-1}x_{\mathrm{sep}}^{3/2},\\
  \lVert C_p\rVert_{C_\eta^k}
    &\leq C_k\mathsf C^{-2}\qquad(x=x_{\mathrm{sep}}).
\end{aligned}
\tag{B.39}\label{eq:B.39}
\end{equation}
The corresponding integrals of the ideal profile with positive weights also vanish as positive
powers of \NSi{NS-F-P156-027}{x_{\mathrm{sep}}}. The pressure increment of the reference profile at
\NSi{NS-F-P156-028}{x_{\mathrm{sep}}} is
\NSi{NS-F-P156-029}{(5/2)P_*^2f^2x_{\mathrm{sep}}^{1/5}=O(\mathsf C^{-2})}. All constants in these vanishing errors may depend on the already fixed axis and
shape parameters.

After \NSi{NS-F-P156-030}{X_{\mathrm{sep}}}, formula \eqref{eq:B.34} is exactly the ideal angular profile because
\NSi{NS-F-P156-031}{\mathsf C^{-1}f(X/X_i)^{1/10}=P_*fx^{1/10}}. Choose the amplitude so large that
\NSi{NS-F-P156-032}{x_{\mathrm{sep}}<e^{-8}}. On \NSi{NS-F-P156-033}{-8<\log x<-7}, restore
\NSi{NS-F-P156-034}{G_i} smoothly to \NSi{NS-F-P156-035}{4\eta}, keeping
\NSi{NS-F-P156-036}{E} ideal. The discrepancy \NSi{NS-F-P156-037}{G_i-4\eta} can be made as small as required
by first choosing \NSi{NS-F-P156-038}{j_0}, then \NSi{NS-F-P156-039}{\Lambda}, and finally the small transitions in \eqref{eq:B.32}. Integration over the
remaining interval gives bounds depending only on the exterior-profile data: all relevant
weights are integrable positive powers at zero. Therefore the five normalized moment
discrepancies at the moment correction have size at most
\NSi{NS-F-P156-040}{C_k\lVert G_i-4\eta\rVert_{C_\eta^{k+1}}+o_{\mathrm C}(1)}. The same estimate controls field values, logarithmic radial derivatives, and
\NSi{NS-F-P156-041}{Q_s,N_s} on the fixed restoration patch.

On \NSi{NS-F-P156-042}{-6<\log x<-5}, add two fixed radial bumps to \NSi{NS-F-P156-043}{U} and three to
\NSi{NS-F-P156-044}{E}. At the ideal profile, replace the \NSi{NS-F-P156-045}{J} row by
\NSi{NS-F-P156-046}{J-4\eta I}, and the \NSi{NS-F-P156-047}{S} row by
\NSi{NS-F-P156-048}{S-8\eta M}. Up to fixed nonzero normalizing factors, the derivative of the normalized moment map has the two blocks
\NSFormula{NS-F-P156-049}{display}{none}
\[
  U\ \text{weights}:\quad 1,\ fx^{3/5};
  \qquad
  E\ \text{weights}:\quad x^{1/2},\ fx^{1/10},\ fx^{-9/10}.
\]
